%----------------------------------------------------------------------------------------
% RATIONALE FUNCTIES EN ASYMPTOTEN. Alle opgaven zijn eigen werk.
%
% Opbouw: Probleem, Start, leerstof (theorie, werkwijze of uitleg, voorbeeld, korte oefening), en op een
% nieuwe pagina de Opdrachten (door elkaar, stijgende moeilijkheid), Koppel en Test jezelf.
% Labels van dit hoofdstuk beginnen met h2-.
%----------------------------------------------------------------------------------------

\mychapter{2}{Rationale functies\alleenvoor{2,3}{ en wortelfuncties}}

\begin{probleem}
\naast{Een sirene zendt een toon van $700$ Hz uit en komt met een snelheid van $v$ m/s op je af. Voor de sirene uit liggen de geluidsgolven dichter op elkaar, en je hoort een hogere toon:
\[F(v) = \frac{700 \cdot 343}{343 - v} \ \text{Hz} ,\]
met $343$ m/s de snelheid van het geluid in lucht. Rijdt de sirene van je weg, dan is $v$ negatief en hoor je een lagere toon. Een ambulance aan $30$ m/s klinkt zo eerst als $767$ Hz en daarna als $644$ Hz. Wat hoor je als de sirene steeds sneller op je afkomt, tot bijna de snelheid van het geluid? En als ze steeds sneller van je weg gaat?}{\begin{tikzpicture}[x=.45cm,y=.45cm,font=\scriptsize]
  \foreach \k in {0,1,2,3}{\draw[accent!70] ({.55*\k},0) circle ({4-\k});}
  \fill[accent] (2.2,0) circle (3.5pt);
  \draw[->,thick,accent] (2.2,0) -- (3.4,0) node[above,text=accent] {$v$};
  \fill (6.3,0) circle (2.5pt) node[below,align=center] {jij};
  \node[below,align=center] at (-2.4,-4.1) {lager};
  \node[below,align=center] at (5.2,-4.1) {hoger};
\end{tikzpicture}}
\end{probleem}

\begin{alleenvoor}{4}
\lkr{Leerplan \pakketnaam: LPD 24-25 (rationale functies), het deel over domein, nulwaarden, perforaties, asymptoten en het schetsen van de grafiek. Limieten komen in de cursus Limieten en afgeleiden; hier worden asymptoten informeel ingevoerd (\textit{de functiewaarden worden onbeperkt groot}, \textit{het verschil met de rechte wordt willekeurig klein}) en bepaald met de euclidische deling uit hoofdstuk 1.}
\end{alleenvoor}
\begin{alleenvoor}{2,3}
\lkr{Leerplan \pakketnaam: LPD \alleenvoor{2}{17}\alleenvoor{3}{18} (het verband tussen de grafiek en de kenmerken van rationale en irrationale functies: domein, nulwaarden, tekenverloop, asymptotisch gedrag en gedrag op oneindig; bij de homografische functies en de $n$-de machtswortelfuncties, de voorbeelden uit de wenk, ook bereik en stijgen en dalen) en LPD 6 (vergelijkingen en ongelijkheden grafisch oplossen met ICT, sectie~\ref{sec:h2-ict}, op de TI-84 met \texttt{intersect} in het menu CALC). Ook de cursus Afgeleiden gebruikt in dit pakket alleen een informeel limietbegrip (wenk bij LPD \alleenvoor{2}{18}\alleenvoor{3}{19}); hier worden asymptoten informeel ingevoerd (\textit{de functiewaarden worden onbeperkt groot}, \textit{het verschil met de rechte wordt willekeurig klein}) en bepaald met de euclidische deling uit hoofdstuk 1. Rationale vergelijkingen en ongelijkheden algebraïsch oplossen (hoofdstuk 3 van pakket 4) vraagt het leerplan niet: hier lossen leerlingen ze grafisch op. De bewijzen zijn \alleenvoor{2}{extra bij LPD 4}\alleenvoor{3}{bewijsvoering bij LPD 4}.}
\end{alleenvoor}

De toon $F(v)$ uit het probleem is een quotiënt van twee veeltermen in $v$. Zulke functies heten rationaal. Invullen lukt: $F(300) \approx 5584$ Hz, en $F(340) \approx 80\,033$ Hz, ver boven wat je oor nog hoort (tot ongeveer $20\,000$ Hz). Maar bij $v = 343$ is de noemer nul, en zegt de formule niets meer. Wat doet de grafiek in de buurt van $343$? En wat gebeurt er als $v$ zeer groot en negatief wordt? Daarvoor heb je asymptoten nodig: rechten waar de grafiek steeds dichter bij komt, zonder ze te bereiken.

Alles uit hoofdstuk 1 komt hier terug. Met ontbinden vind je het domein, de nulwaarden en de gaten in de grafiek. Met de euclidische deling vind je hoe de grafiek zich gedraagt voor grote $|x|$. En de tekentabel zegt aan welke kant van de asymptoten de grafiek ligt.\alleenvoor{2,3}{ Op het einde bestudeer je de wortelfuncties, en los je vergelijkingen en ongelijkheden met breuken en wortels grafisch op met ICT.}

\begin{oefeningen}[Start]
\oef{h2-start-1}

\oef{h2-verken}
\end{oefeningen}

%========================================================================================
\section{Domein, nulwaarden en perforaties}

\subsection{Domein en nulwaarden}

Voor welke snelheden kun je de formule uit het probleem niet gebruiken? Alleen voor $v = 343$: daar is de noemer nul. Een breuk bestaat alleen als de noemer niet nul is. Het domein van een rationale functie vind je dus door de nulwaarden van de noemer te zoeken, en dat doe je door te ontbinden.

\begin{theorie}
Een \textbf{\textcolor{accent}{rationale functie}} is een quotiënt van twee veeltermfuncties: \quad $\formule{f(x) = \frac{A(x)}{B(x)}}$

Het \textbf{\textcolor{accent}{domein}} bestaat uit alle reële getallen behalve de nulwaarden van de noemer.

\begin{stelling}{Nulwaarden van een rationale functie}
$a$ is een nulwaarde van $f(x) = \frac{A(x)}{B(x)}$ \quad $\formule{\iff}$ \quad $\formule{A(a) = 0 \text{ en } B(a) \neq 0}$
\end{stelling}

\begin{bewijs}
Een nulwaarde moet tot het domein behoren, dus $B(a) \neq 0$. Voor zo'n $a$ geldt
\bewijsregel{\frac{A(a)}{B(a)} = 0 \iff A(a) = 0 ,}
want je mag beide leden vermenigvuldigen met het getal $B(a) \neq 0$, en terug delen.
\end{bewijs}
\end{theorie}

Een nulwaarde van de teller is dus niet altijd een nulwaarde van de functie: ze moet ook in het domein liggen. Bepaal daarom altijd eerst het domein.

\voorbeeld{}{Bepaal het domein en de nulwaarden van $f(x) = \dfrac{x^3 - 4x}{x^2 + x - 6}$.\label{vb:h2-nul}}{%
\uitleg{Ontbind teller en noemer. Uit de noemer lees je het domein af, uit de teller de kandidaat-nulwaarden.}
\begin{stappenbreed}
f(x) = \frac{x(x - 2)(x + 2)}{(x - 2)(x + 3)} & Teller: zonder $x$ af en $x^2 - 4$. Noemer: som $-1$, product $-6$. \\
\dom f = \R \setminus \{-3, 2\} & De nulwaarden van de noemer. \\
\text{teller nul in } 0,\ 2,\ -2 & Kandidaten. \\
\text{nulwaarden: } -2 \text{ en } 0 & $2$ ligt niet in het domein.
\end{stappenbreed}}

\begin{oefeningen}
\needspace{6\baselineskip}
\oef{h2-nul}
\end{oefeningen}

\subsection{Perforaties}

Wat gebeurt er in voorbeeld~\ref{vb:h2-nul} bij $x = 2$? Teller en noemer hebben daar allebei de factor $x - 2$. Voor $x \neq 2$ mag je die schrappen, want dan deel je teller en noemer door hetzelfde getal $x - 2 \neq 0$. De grafiek is dan die van een eenvoudiger functie, met één punt eruit.

\begin{theorie}
Is $f(x) = \dfrac{(x - a) \cdot A_1(x)}{(x - a) \cdot B_1(x)}$, dan is $\formule{f(x) = \frac{A_1(x)}{B_1(x)}}$ voor $x \neq a$.

Is $B_1(a) \neq 0$, dan heeft de grafiek van $f$ een \textbf{\textcolor{accent}{perforatie}} in $\formule{\Bigl(a, \frac{A_1(a)}{B_1(a)}\Bigr)}$: een gat in de grafiek van de vereenvoudigde functie.
\end{theorie}

\voorbeeld{}{Bepaal domein, nulwaarden en perforaties van $f(x) = \dfrac{x^2 - 4}{x^2 - x - 2}$.\label{vb:h2-perforatie}}{%
\naast{%
\uitleg{Ontbind teller en noemer. Het domein lees je af vóór je vereenvoudigt; de nulwaarden en de perforatie erna.}
\begin{stappenvrij}{4.4cm}{5.3cm}
f(x) = \frac{(x - 2)(x + 2)}{(x - 2)(x + 1)} & Ontbinden. \\
\dom f = \R \setminus \{-1, 2\} & Nulwaarden van de noemer. \\
f(x) = \frac{x + 2}{x + 1} \text{ voor } x \neq 2 & Factor $x - 2$ schrappen. \\
\text{nulwaarde } -2 & De teller van wat overblijft. \\
\text{perforatie in } \bigl(2, \tfrac43\bigr) & $\frac{2 + 2}{2 + 1} = \frac43$.
\end{stappenvrij}}{%
\begin{tikzpicture}[x=.42cm,y=.42cm,font=\scriptsize]
  \rooster{-5}{1}{4}{-3}{1}{5}
  \kromme{-5}{-1.25}{(\x+2)/(\x+1)} \kromme{-.75}{4}{(\x+2)/(\x+1)}
  \vol{(-2,0)} \open{(2,4/3)}
\end{tikzpicture}}}

In een perforatie is $f$ niet gedefinieerd, maar de grafiek loopt er vlak langs: links en rechts van het gat liggen de punten van $\frac{A_1}{B_1}$. Je tekent het gat als een open bolletje.

\begin{oefeningen}
\needspace{6\baselineskip}
\oef{h2-domein}
\end{oefeningen}

%========================================================================================
\section{Asymptoten}

\subsection{Verticale asymptoten}

Wat hoor je in het probleem als $v$ dicht bij $343$ komt? Reken maar uit:
\[F(340) \approx 80\,033, \qquad F(342{,}9) = 2\,401\,000, \qquad F(342{,}99) = 24\,010\,000 .\]
Hoe dichter $v$ bij $343$, hoe groter $F(v)$, zonder grens. Hetzelfde zie je bij de eenvoudigste rationale functie.

\naast{%
Bij $f(x) = \frac{1}{x}$ worden de functiewaarden onbeperkt groot als $x$ dicht bij $0$ komt: $f(0{,}001) = 1000$ en $f(-0{,}001) = -1000$. De grafiek schiet langs de $y$-as omhoog en omlaag, zonder die ooit te raken. Zo'n rechte heet een verticale asymptoot.

Een perforatie ontstaat als de noemer nul wordt door een factor die je kunt schrappen. Blijft de factor na het schrappen in de noemer staan, dan krijg je een asymptoot.

\textit{Dicht bij}, \textit{onbeperkt groot} en \textit{willekeurig klein} gebruiken we hier informeel. Met limieten worden de definities en de bewijzen van dit hoofdstuk later sluitend.}{%
\begin{tikzpicture}[x=.42cm,y=.42cm,font=\scriptsize]
  \rooster{-4}{1}{4}{-4}{1}{4}
  \kromme{.25}{4}{1/\x} \kromme{-4}{-.25}{1/\x}
  \node[accent,above right] at (.5,2) {$\frac1x$};
\end{tikzpicture}}

\begin{theorie}
De rechte $x = a$ is een \textbf{\textcolor{accent}{verticale asymptoot}} van $f$ als $|f(x)|$ onbeperkt groot wordt als $x$ dicht bij $a$ komt.

\begin{stelling}{Verticale asymptoot}
Hebben $A(x)$ en $B(x)$ geen gemeenschappelijke factor, dan geeft elke nulwaarde $a$ van de noemer $B(x)$ een verticale asymptoot \quad $\formule{x = a}$
\end{stelling}

\begin{bewijs}
Was $A(a) = 0$, dan was $x - a$ volgens de factorstelling een factor van $A(x)$ en van $B(x)$. Dus is $A(a) \neq 0$. Komt $x$ dicht bij $a$, dan komt $A(x)$ dicht bij $A(a) \neq 0$, want de grafiek van een veeltermfunctie maakt geen sprongen, en $B(x)$ komt dicht bij $B(a) = 0$. Een getal dicht bij $A(a)$ gedeeld door een getal dicht bij $0$ heeft een grote absolute waarde: zo is $\frac{2}{0{,}001} = 2000$. Dus wordt
\bewijsregel{|f(x)| = \frac{|A(x)|}{|B(x)|}}
onbeperkt groot.
\end{bewijs}
\end{theorie}

Een nulwaarde van de vereenvoudigde noemer heet ook een \textbf{pool}. Let op de voorwaarde: teller en noemer mogen geen gemeenschappelijke factor hebben. Schrap die dus eerst, en schrap ze helemaal: soms komt een factor meer dan eens voor.

\voorbeeld{}{Heeft $f(x) = \dfrac{x^2 - 2x}{x^2 - 4x + 4}$ in $2$ een perforatie of een verticale asymptoot?\label{vb:h2-ofva}}{%
\uitleg{Ontbind en schrap wat je kunt schrappen. Kijk dan of $x - 2$ nog in de noemer staat.}
\begin{stappenbreed}
f(x) = \frac{x(x - 2)}{(x - 2)^2} & De noemer is een kwadraat. \\[2ex]
f(x) = \frac{x}{x - 2} \text{ voor } x \neq 2 & Eén factor $x - 2$ valt weg, er blijft er één over. \\[1ex]
\text{verticale asymptoot } x = 2 & De noemer is nog nul in $2$, de teller niet.
\end{stappenbreed}
Wie na één keer schrappen besluit dat er een perforatie is, vergist zich: in $2$ is er geen functiewaarde om in te vullen.}

\begin{werkwijze}
Ontbind teller en noemer, en schrap alle gemeenschappelijke factoren. Neem elke nulwaarde $a$ van de oorspronkelijke noemer. Staat $x - a$ nog in de vereenvoudigde noemer, dan is $x = a$ een verticale asymptoot. Is $x - a$ helemaal weggevallen, dan is er een perforatie: vul $a$ in de vereenvoudigde functie in.
\end{werkwijze}

\begin{oefeningen}
\needspace{6\baselineskip}
\oef{h2-va}
\end{oefeningen}

Een verticale asymptoot zegt dat de grafiek onbeperkt stijgt of daalt, maar niet in welke richting. Dat lees je af in de tekentabel. Is $f$ net links van de pool negatief, dan gaat de grafiek daar naar onder; is $f$ positief, naar boven. Voor de noemer geldt dezelfde regel als voor een veelterm in hoofdstuk 1: bij een pool met oneven multipliciteit wisselt het teken, bij een even multipliciteit niet. De teller verandert er niet van teken, want ze is niet nul in de pool.

\voorbeeld{}{Bepaal de verticale asymptoten van $f(x) = \dfrac{x - 3}{x^3 - 4x^2 + 4x}$. Gaat de grafiek links en rechts ervan naar boven of naar onder?\label{vb:h2-vazijde}}{%
\naast{%
\uitleg{Ontbind de noemer. Stel de tekentabel op, met de nulwaarde van de teller en de polen.}
\begin{stappenvrij}{4.3cm}{5.6cm}
f(x) = \frac{x - 3}{x(x - 2)^2} & Geen gemeenschappelijke factor. \\[1ex]
x = 0,\quad x = 2 & Verticale asymptoten. \\[1ex]
x \cdot x^2 \text{ in de noemer} & Rechts van $3$: $f > 0$. \\
2:\ \text{dubbele pool} & Teken blijft. \\
0:\ \text{enkele pool} & Teken wisselt.
\end{stappenvrij}
\figuren{\tekentabel{cccccccc}{x & & 0 & & 2 & & 3 & }{f(x) & + & | & - & | & - & 0 & +}}
Bij $x = 0$: links naar boven, rechts naar onder. Bij $x = 2$: aan beide kanten naar onder.}{%
\begin{tikzpicture}[x=.42cm,y=.42cm,font=\scriptsize]
  \rooster{-3}{1}{5}{-4}{1}{4}
  \draw[grijs,dashed] (2,-4) -- (2,4.3);
  \kromme{-3}{-.17}{(\x-3)/(\x*(\x-2)^2)}
  \kromme{.22}{1.5}{(\x-3)/(\x*(\x-2)^2)}
  \kromme{2.28}{5}{(\x-3)/(\x*(\x-2)^2)}
  \vol{(3,0)}
\end{tikzpicture}}}

\begin{oefeningen}
\needspace{6\baselineskip}
\oef{h2-tekentabel}
\end{oefeningen}

\subsection{Horizontale asymptoot: teller van lagere graad}

En als de sirene steeds sneller van je weg gaat? Dan is $v$ negatief, en
\[F(-1000) \approx 179, \qquad F(-10\,000) \approx 23, \qquad F(-100\,000) \approx 2{,}4 .\]
De toon zakt naar $0$ Hz, maar bereikt die nooit, want de teller is nooit nul. Voor grote $|x|$ kijk je dus naar de graden van teller en noemer. In de Verken kwam $G(q) = \frac{500 + 2q}{q}$ steeds dichter bij $2$: de grafiek nadert de rechte $y = 2$. We beginnen met het geval waarin de noemer de hoogste graad heeft.

\begin{theorie}
De rechte $y = c$ is een \textbf{\textcolor{accent}{horizontale asymptoot}} van $f$ als $f(x)$ willekeurig dicht bij $c$ komt voor grote $|x|$.

\begin{stelling}{Horizontale asymptoot $y = 0$}
Is graad $A <$ graad $B$, dan heeft $f(x) = \frac{A(x)}{B(x)}$ de horizontale asymptoot \quad $\formule{y = 0}$
\end{stelling}

\begin{bewijs}
Stel dat $A$ graad $n$ heeft, en $B$ graad $m > n$ met hoogstegraadscoëfficiënt $b_m$. Deel teller en noemer door $x^m$:
\bewijsregel{f(x) = \frac{\dfrac{a_n}{x^{m-n}} + \dfrac{a_{n-1}}{x^{m-n+1}} + \cdots + \dfrac{a_0}{x^m}}{b_m + \dfrac{b_{m-1}}{x} + \cdots + \dfrac{b_0}{x^m}} .}
Elke breuk is een vast getal gedeeld door een macht van $x$, en wordt willekeurig klein voor grote $|x|$. In de teller staan alleen zulke breuken, dus de teller wordt willekeurig klein. De noemer komt dicht bij $b_m \neq 0$. Een heel klein getal gedeeld door een getal dicht bij $b_m$ is heel klein: $f(x)$ komt willekeurig dicht bij $0$.
\end{bewijs}
\end{theorie}

Een grafiek kan haar horizontale asymptoot wel snijden, in tegenstelling tot een verticale. Een asymptoot zegt alleen iets over grote $|x|$. Ze snijdt de asymptoot $y = 0$ precies in de nulwaarden van $f$.

\voorbeeld{}{Bepaal de asymptoten van $f(x) = \dfrac{3x + 6}{x^2 + 2}$. Snijdt de grafiek de horizontale asymptoot?\label{vb:h2-ha0}}{%
\naast{%
\uitleg{De noemer is nooit nul, en de teller heeft een lagere graad dan de noemer.}
\begin{stappenvrij}{4.6cm}{4.5cm}
x^2 + 2 > 0 & Geen verticale asymptoot. \\[1ex]
y = 0 & Graad $1 <$ graad $2$. \\[1ex]
3x + 6 = 0 \iff x = -2 & Snijpunt $(-2, 0)$.
\end{stappenvrij}
Voor $x < -2$ ligt de grafiek onder de asymptoot, en ze komt er van onderen naar terug.}{%
\begin{tikzpicture}[x=.38cm,y=.5cm,font=\scriptsize]
  \rooster{-6}{1}{6}{-1}{1}{4}
  \kromme{-6}{6}{(3*\x+6)/((\x)^2+2)}
  \vol{(-2,0)}
\end{tikzpicture}}}

\begin{oefeningen}
\needspace{6\baselineskip}
\oef{h2-ha0}
\end{oefeningen}

\needspace{26\baselineskip}
Nu kun je het probleem van de sirene oplossen.

\voorbeeld{}{Terug naar het probleem. Een sirene van $700$ Hz komt met snelheid $v$ (in m/s) op je af. Je hoort dan de toon $F(v) = \dfrac{240\,100}{343 - v}$ Hz. Bepaal de asymptoten van de grafiek van $F$. Wat betekenen ze?\label{vb:h2-terug}}{%
\uitleg{Hier is $240\,100 = 700 \cdot 343$. De teller is een getal, de noemer heeft graad 1: er is geen gemeenschappelijke factor. Het teken van $F$ hangt alleen van de noemer af.}
\naast{\begin{stappenvrij}{4.1cm}{4.3cm}
\dom F = \R \setminus \{343\} & De noemer is nul in $343$. \\[1ex]
v = 343 & Verticale asymptoot: de teller is nooit nul. \\[1ex]
F(v) > 0 \text{ als } v < 343 & Links van de asymptoot gaat de grafiek naar boven. \\[1ex]
F = 0 & Horizontale asymptoot: graad $0 <$ graad $1$. Voor grote negatieve $v$ ligt de grafiek erboven.
\end{stappenvrij}}{%
\begin{tikzpicture}[x=.36cm,y=.62cm,font=\scriptsize]
  \draw[black!14,very thin] (-6,0) grid (4,4);
  \draw[->,grijs] (-6,0) -- (4.4,0) node[right,text=black] {$v$};
  \draw[->,grijs] (0,0) -- (0,4.4) node[above,text=black] {$F$};
  \foreach \t in {-6,-4,-2,2}{\node[below,inner sep=1.5pt] at (\t,0) {$\the\numexpr\t*100\relax$};}
  \foreach \t in {1,2,3,4}{\node[left,inner sep=1.5pt] at (0,\t) {$\t\,000$};}
  \draw[grijs,dashed] (3.43,0) -- (3.43,4.4) node[above,text=black] {$v = 343$};
  \kromme{-6}{2.83}{2.401/(3.43-\x)}
  \vol{(0,.7)}
\end{tikzpicture}}
Hoe dichter de sirene bij de snelheid van het geluid komt, hoe hoger de toon, zonder grens. In werkelijkheid stapelen de golven zich dan voor de bron op tot een schokgolf: de geluidsmuur. Hoe sneller de sirene van je weg gaat, hoe lager de toon, maar $0$ Hz wordt het nooit. Voor ${v > 343}$ is $F(v)$ negatief: de formule heeft daar geen betekenis, want de bron is sneller dan haar eigen geluid, en je hoort haar pas als ze al voorbij is.}

\subsection{Horizontale asymptoot: teller en noemer van dezelfde graad}

Hebben teller en noemer dezelfde graad, dan helpt de euclidische deling. Het quotiënt is dan een getal, en de rest heeft een lagere graad dan de noemer: dat geval ken je al.

\begin{theorie}
\begin{stelling}{Horizontale asymptoot $y = c$}
Is graad $A =$ graad $B$, met hoogstegraadscoëfficiënten $a_n$ en $b_n$, dan heeft $f(x) = \frac{A(x)}{B(x)}$ de horizontale asymptoot \quad $\formule{y = \frac{a_n}{b_n}}$
\end{stelling}

\begin{bewijs}
Neem $c = \frac{a_n}{b_n}$. In $R(x) = A(x) - c \cdot B(x)$ vallen de termen in $x^n$ weg, want $c \cdot b_n = a_n$. Dus is $R = 0$ of graad $R < n$, en
\bewijsregel{f(x) = \frac{c \cdot B(x) + R(x)}{B(x)} = c + \frac{R(x)}{B(x)} .}
Volgens de vorige stelling wordt $\frac{R(x)}{B(x)}$ willekeurig klein voor grote $|x|$, dus komt $f(x)$ willekeurig dicht bij $c$.
\end{bewijs}
\end{theorie}

Het bewijs geeft meer dan de asymptoot: het teken van $\frac{R(x)}{B(x)}$ zegt of de grafiek boven of onder de asymptoot ligt.

\voorbeeld{}{Bepaal de asymptoten van $f(x) = \dfrac{2x^2 + 1}{x^2 - 1}$. Ligt de grafiek boven of onder de horizontale asymptoot?\label{vb:h2-ha}}{%
\naast{%
\uitleg{De teller is nooit nul, dus er is geen gemeenschappelijke factor. Teller en noemer hebben dezelfde graad.}
\begin{stappenvrij}{4.5cm}{4.9cm}
x^2 - 1 = 0 \iff x = \pm 1 & Verticale asymptoten $x = -1$ en $x = 1$. \\[1ex]
y = \frac{2}{1} = 2 & Horizontale asymptoot. \\[1ex]
f(x) = 2 + \frac{3}{x^2 - 1} & $2x^2 + 1 = 2(x^2 - 1) + 3$. \\[1ex]
\frac{3}{x^2 - 1} > 0 \text{ voor } |x| > 1 & Daar boven de asymptoot.
\end{stappenvrij}}{%
\begin{tikzpicture}[x=.42cm,y=.36cm,font=\scriptsize]
  \rooster{-4}{1}{4}{-4}{1}{7}
  \draw[grijs,dashed] (-1,-4) -- (-1,7.3) (1,-4) -- (1,7.3) (-4,2) -- (4.3,2);
  \kromme{-4}{-1.26}{2+3/((\x)^2-1)} \kromme{-.71}{.71}{2+3/((\x)^2-1)} \kromme{1.26}{4}{2+3/((\x)^2-1)}
\end{tikzpicture}}}

\begin{oefeningen}
\needspace{6\baselineskip}
\oef{h2-ha}
\end{oefeningen}

\subsection{Schuine asymptoot}

Is de graad van de teller één meer dan die van de noemer, dan is het quotiënt van de deling een eerstegraadsveelterm. De grafiek nadert dan geen horizontale maar een schuine rechte.

\begin{theorie}
De rechte $y = mx + q$ met $m \neq 0$ is een \textbf{\textcolor{accent}{schuine asymptoot}} van $f$ als het verschil $f(x) - (mx + q)$ willekeurig klein wordt voor grote $|x|$.

\begin{stelling}{Schuine asymptoot}
Is graad $A =$ graad $B + 1$ en is $Q(x)$ het quotiënt van de deling $A(x) : B(x)$, dan heeft $f(x) = \frac{A(x)}{B(x)}$ de schuine asymptoot \quad $\formule{y = Q(x)}$
\end{stelling}

\begin{bewijs}
Uit de euclidische deling $A = B \cdot Q + R$ volgt
\bewijsregel{f(x) = Q(x) + \frac{R(x)}{B(x)} \qquad \text{met } R = 0 \text{ of graad } R < \text{graad } B .}
De rest heeft een lagere graad dan $B$, en dus ook dan $A$. Daarom heeft $B \cdot Q$ dezelfde graad als $A$, één meer dan $B$, en is graad $Q = 1$ (graad van een product): $Q(x) = mx + q$ met $m \neq 0$. Het verschil $f(x) - Q(x) = \frac{R(x)}{B(x)}$ wordt willekeurig klein voor grote $|x|$, volgens de stelling over de asymptoot $y = 0$.
\end{bewijs}
\end{theorie}

\voorbeeld{}{Bepaal de asymptoten van $f(x) = \dfrac{x^2 + x - 1}{x - 1}$.\label{vb:h2-sa}}{%
\naast{%
\uitleg{De graad van de teller is één meer dan die van de noemer: deel met Horner.}
\begin{stappenvrij}{4.2cm}{5.8cm}
f(x) = x + 2 + \frac{1}{x - 1} & Horner met $a = 1$ op $1, 1, -1$ geeft $1, 2$ en rest $1$. \\[1ex]
x = 1 & Verticale asymptoot: de teller is $1 \neq 0$ in $1$. \\[1ex]
y = x + 2 & Schuine asymptoot. \\[1ex]
\frac{1}{x - 1} > 0 \text{ voor } x > 1 & Rechts boven de asymptoot, links eronder.
\end{stappenvrij}}{%
\begin{tikzpicture}[x=.4cm,y=.3cm,font=\scriptsize]
  \rooster{-4}{1}{5}{-4}{2}{8}
  \draw[grijs,dashed] (1,-4) -- (1,8.5) (-4,-2) -- (5.3,7.3);
  \kromme{-4}{.85}{\x+2+1/(\x-1)} \kromme{1.17}{5}{\x+2+1/(\x-1)}
\end{tikzpicture}}}

\begin{oefeningen}
\needspace{6\baselineskip}
\oef{h2-sa}
\end{oefeningen}

\subsection{Teller van een veel hogere graad}

Is de graad van de teller twee of meer hoger dan die van de noemer, dan heeft het quotiënt $Q(x)$ graad $2$ of meer. Het bewijs van de vorige stelling werkt nog altijd: het verschil $f(x) - Q(x)$ wordt willekeurig klein. Maar de grafiek van $Q$ is nu geen rechte meer. Er is dan geen horizontale of schuine asymptoot, en de grafiek gedraagt zich voor grote $|x|$ zoals die van $Q$.

\voorbeeld{}{Heeft de grafiek van $f(x) = \dfrac{x^3 - 2}{x}$ een horizontale of een schuine asymptoot?\label{vb:h2-geen}}{%
\naast{%
\uitleg{De teller heeft graad 3, de noemer graad 1. Splits de breuk.}
\begin{stappenvrij}{4.4cm}{5.4cm}
f(x) = \frac{x^3}{x} - \frac{2}{x} = x^2 - \frac{2}{x} & Het quotiënt is $x^2$, de rest $-2$. \\[1ex]
\text{geen} & Het quotiënt heeft graad 2. \\[1ex]
x = 0 & Wel een verticale asymptoot.
\end{stappenvrij}
Voor grote $|x|$ ligt de grafiek dicht bij de parabool $y = x^2$ (gestippeld).}{%
\begin{tikzpicture}[x=.5cm,y=.27cm,font=\scriptsize]
  \rooster{-3}{1}{3}{-4}{2}{8}
  \draw[grijs,dashed] plot[domain=-2.9:2.9,samples=60,smooth] (\x,{(\x)^2});
  \kromme{-2.88}{-.22}{(\x)^2-2/\x} \kromme{.47}{2.95}{(\x)^2-2/\x}
\end{tikzpicture}}}

Samengevat beslist de graad van de teller tegenover die van de noemer welke asymptoot er is voor grote $|x|$. Je vindt ze altijd met de euclidische deling $f(x) = Q(x) + \frac{R(x)}{B(x)}$: de asymptoot is $y = Q(x)$, zolang die een rechte is.

\begin{center}
\renewcommand{\arraystretch}{2.1}
\begin{tabular}{|c|c|c|c|}
\hline
\textbf{graden} & \textbf{quotiënt} $Q(x)$ & \textbf{asymptoot voor grote} $|x|$ & \textbf{voorbeeld} \\ \hline
graad $A <$ graad $B$ & $0$ & horizontaal: $y = 0$ & $\dfrac{3x + 6}{x^2 + 2}$ \\ \hline
graad $A =$ graad $B$ & $\dfrac{a_n}{b_n}$ & horizontaal: $y = \dfrac{a_n}{b_n}$ & $\dfrac{2x^2 + 1}{x^2 - 1}$ \\ \hline
graad $A =$ graad $B + 1$ & $mx + q$ & schuin: $y = mx + q$ & $\dfrac{x^2 + x - 1}{x - 1}$ \\ \hline
graad $A \geqslant$ graad $B + 2$ & graad $\geqslant 2$ & geen & $\dfrac{x^3 - 2}{x}$ \\ \hline
\end{tabular}
\end{center}

\begin{oefeningen}
\needspace{6\baselineskip}
\oef{h2-soort}

\oef{h2-asymptoten}
\end{oefeningen}

%========================================================================================
\section{De grafiek schetsen}

\subsection{De zeven stappen}

Je kent nu alle bouwstenen. Om de grafiek te schetsen, verzamel je ze in een vaste volgorde. De tekentabel bindt alles samen: ze zegt waar de grafiek boven of onder de $x$-as ligt, ook naast de asymptoten.

\voorbeeld{}{Schets de grafiek van $f(x) = \dfrac{x - 1}{x^2 - 4}$.\label{vb:h2-schets}}{%
\naast{%
\uitleg{Verzamel alle informatie en zet ze in een tekentabel; die zegt waar de grafiek boven of onder de $x$-as ligt, ook naast de asymptoten.}
\begin{stappenvrij}{4cm}{3.3cm}
\dom f = \R \setminus \{-2, 2\} & Geen perforaties. \\
\text{nulwaarde } 1 & \\
x = -2,\ x = 2 & Verticale asymptoten. \\
y = 0 & Graad teller $<$ graad noemer. \\
f(0) = \tfrac14 & Snijpunt met de $y$-as.
\end{stappenvrij}
\figuren{\tekentabel{cccccccc}{x & & -2 & & 1 & & 2 & }{f(x) & - & | & + & 0 & - & | & +}}}{%
\begin{tikzpicture}[x=.48cm,y=.48cm,font=\scriptsize]
  \rooster{-5}{1}{5}{-3}{1}{3}
  \draw[grijs,dashed] (-2,-3) -- (-2,3.3) (2,-3) -- (2,3.3);
  \kromme{-5}{-2.18}{(\x-1)/((\x)^2-4)}
  \kromme{-1.83}{1.87}{(\x-1)/((\x)^2-4)}
  \kromme{2.06}{5}{(\x-1)/((\x)^2-4)}
  \vol{(1,0)}
\end{tikzpicture}}}

\begin{werkwijze}
\textbf{1.} Domein. \textbf{2.} Ontbind en vereenvoudig: perforaties. \textbf{3.} Nulwaarden. \textbf{4.} Verticale asymptoten. \textbf{5.} Horizontale of schuine asymptoot. \textbf{6.} Tekentabel met nulwaarden en polen. \textbf{7.} Een paar functiewaarden, en schetsen: de grafiek nadert de asymptoten aan de kant die de tekentabel aangeeft.
\end{werkwijze}

\begin{oefeningen}
\oef{h2-schets}
\end{oefeningen}

In het volgende voorbeeld komt alles samen: een perforatie, een verticale en een schuine asymptoot.

\voorbeeld{}{Schets de grafiek van $f(x) = \dfrac{x^3 - x}{x^2 - x - 2}$.\label{vb:h2-schets2}}{%
\naast{%
\uitleg{Volg de zeven stappen. Werk na het schrappen verder met de vereenvoudigde functie.}
\begin{stappenvrij}{5.3cm}{4.7cm}
f(x) = \frac{x(x - 1)(x + 1)}{(x - 2)(x + 1)} & $\dom f = \R \setminus \{-1, 2\}$. \\[2ex]
f(x) = \frac{x(x - 1)}{x - 2} \text{ voor } x \neq -1 & Perforatie in $\bigl(-1, -\tfrac23\bigr)$. \\[2ex]
\text{nulwaarden } 0,\ 1 & \\
x = 2 & Verticale asymptoot. \\
f(x) = x + 1 + \frac{2}{x - 2} & Horner op $1, -1, 0$ met $a = 2$. \\[2ex]
y = x + 1 & Schuine asymptoot.
\end{stappenvrij}
\figuren{\tekentabel{cccccccc}{x & & 0 & & 1 & & 2 & }{f(x) & - & 0 & + & 0 & - & | & +}}}{%
\begin{tikzpicture}[x=.36cm,y=.3cm,font=\scriptsize]
  \rooster{-4}{1}{6}{-5}{2}{9}
  \draw[grijs,dashed] (2,-5) -- (2,9.5) (-4,-3) -- (6.3,7.3);
  \kromme{-4}{1.74}{\x*(\x-1)/(\x-2)} \kromme{2.35}{6}{\x*(\x-1)/(\x-2)}
  \vol{(0,0)} \vol{(1,0)} \open{(-1,-2/3)}
\end{tikzpicture}}}

De tekentabel is die van de vereenvoudigde functie: in de perforatie verandert het teken niet. Teken het gat als een open bolletje op de grafiek.

\begin{oefeningen}
\needspace{6\baselineskip}
\oef{h2-schets2}
\end{oefeningen}

\subsection{Homografische functies}

De eenvoudigste rationale functies hebben een teller en een noemer van graad 1, zoals $\frac{2x + 1}{x - 3}$. Hun grafiek is altijd die van $\frac1x$, verschoven en uitgerekt. Dan heb je de zeven stappen niet nodig: de transformaties van grafieken die je al kent, volstaan.

\begin{theorie}
Een \textbf{\textcolor{accent}{homografische functie}} is een functie $\formule{f(x) = \frac{\alpha x + \beta}{\gamma x + \delta}}$ met $\gamma \neq 0$.

\begin{stelling}{Homografische functie}
Elke homografische functie is te schrijven als

\qquad $\formule{f(x) = c + \frac{a}{x - p}}$ \quad met \quad $\formule{p = -\frac{\delta}{\gamma}}$ \quad en \quad $\formule{c = \frac{\alpha}{\gamma}}$

Is $a \neq 0$, dan ontstaat haar grafiek uit die van $\frac1x$, met de asymptoten $\formule{x = p}$ en $\formule{y = c}$.
\end{stelling}

\begin{bewijs}
Deel de teller door de noemer:
\bewijsregel{\alpha x + \beta = \frac{\alpha}{\gamma}(\gamma x + \delta) + \Bigl(\beta - \frac{\alpha\delta}{\gamma}\Bigr) .}
Deel door $\gamma x + \delta = \gamma(x - p)$, en noem $a = \frac{1}{\gamma}\bigl(\beta - \frac{\alpha\delta}{\gamma}\bigr)$:
\bewijsregel{f(x) = \frac{\alpha}{\gamma} + \frac{\beta - \frac{\alpha\delta}{\gamma}}{\gamma(x - p)} = c + \frac{a}{x - p} .}
De grafiek van $\frac{a}{x - p} + c$ ontstaat uit die van $\frac1x$ door een verticale uitrekking met factor $a$ (met een spiegeling als $a < 0$), een verschuiving over $p$ naar rechts en over $c$ naar boven. De asymptoten $x = 0$ en $y = 0$ van $\frac1x$ verschuiven mee.
\end{bewijs}
\end{theorie}

De formule voor $a$ hoef je niet te kennen: je deelt telkens zelf. Wel moet je weten wanneer het misloopt. Is $a = 0$, dan is de teller een veelvoud van de noemer, zoals in $\frac{2x + 4}{x + 2}$. Dan is $f(x) = 2$ voor $x \neq -2$: een horizontale rechte met een perforatie, zonder asymptoten.

\voorbeeld{}{Schrijf $f(x) = \dfrac{4x - 2}{2x + 2}$ als $c + \frac{a}{x - p}$. Hoe ontstaat de grafiek uit die van $\frac1x$?\label{vb:h2-homo}}{%
\naast{%
\uitleg{Schrijf de teller als een veelvoud van de noemer plus een getal. Let op: de noemer is $2x + 2 = 2(x + 1)$.}
\begin{stappenvrij}{5.6cm}{4.4cm}
4x - 2 = 2(2x + 2) - 6 & $2 \cdot (2x + 2) = 4x + 4$. \\[1ex]
f(x) = 2 - \frac{6}{2x + 2} = 2 - \frac{3}{x + 1} & Deel teller en noemer door $2$. \\[2ex]
a = -3,\ p = -1,\ c = 2 & \\[1ex]
x = -1,\quad y = 2 & Asymptoten.
\end{stappenvrij}
Rek $\frac1x$ verticaal uit met factor $3$ en spiegel om de $x$-as, schuif $1$ naar links en $2$ naar boven.}{%
\begin{tikzpicture}[x=.4cm,y=.4cm,font=\scriptsize]
  \rooster{-5}{1}{4}{-3}{1}{7}
  \draw[grijs,dashed] (-1,-3) -- (-1,7.3) (-5,2) -- (4.3,2);
  \draw[grijs,thick] plot[domain=-5:-.33,samples=60,smooth] (\x,{1/\x}) plot[domain=.14:4,samples=60,smooth] (\x,{1/\x});
  \kromme{-5}{-1.6}{2-3/(\x+1)} \kromme{-.4}{4}{2-3/(\x+1)}
  \node[grijs,above right] at (.3,2.6) {$\frac1x$};
\end{tikzpicture}}}

\begin{alleenvoor}{2,3}
Uit de grafiek lees je ook de andere kenmerken af. De grafiek van $\frac1x$ daalt links en rechts van $0$, en komt nooit op hoogte $0$. Na het uitrekken en verschuiven ontbreekt de waarde $c$ in het bereik, en een negatieve factor $a$ keert stijgen en dalen om.

\begin{theorie}
\begin{stelling}{Kenmerken van een homografische functie}
Voor $f(x) = c + \frac{a}{x - p}$ met $a \neq 0$ is \quad $\formule{\dom f = \R \setminus \{p\}}$ \quad en \quad $\formule{\ber f = \R \setminus \{c\}}$

Is $a > 0$, dan is $f$ \textbf{\textcolor{accent}{dalend}} op $\left]-\infty, p\right[$ en op $\left]p, +\infty\right[$. Is $a < 0$, dan is $f$ \textbf{\textcolor{accent}{stijgend}} op elk van die delen.
\end{stelling}

\begin{bewijs}
\textit{Bereik.} Neem een getal $y \neq c$. Dan is
\bewijsregel{c + \frac{a}{x - p} = y \iff x - p = \frac{a}{y - c} \iff x = p + \frac{a}{y - c} ,}
en dat getal ligt in het domein, want $\frac{a}{y - c} \neq 0$. De waarde $c$ wordt nooit bereikt, want $\frac{a}{x - p} \neq 0$.

\textit{Stijgen en dalen.} Neem $x_1 < x_2$, allebei rechts of allebei links van $p$. Dan hebben $x_1 - p < x_2 - p$ hetzelfde teken, en dan is $\frac{1}{x_1 - p} > \frac{1}{x_2 - p}$, zoals $\frac12 > \frac13$ en $-\frac12 > -1$. Is $a > 0$, dan blijft de ongelijkheid na vermenigvuldigen met $a$ en optellen van $c$: $f(x_1) > f(x_2)$, $f$ daalt. Is $a < 0$, dan keert ze om: $f$ stijgt.
\end{bewijs}
\end{theorie}

In voorbeeld~\ref{vb:h2-homo} is $a = -3 < 0$: $f$ stijgt op $\left]-\infty, -1\right[$ en op $\left]-1, +\infty\right[$, en $\ber f = \R \setminus \{2\}$. Op het hele domein stijgt $f$ niet, want $f(0) = -1 < 5 = f(-2)$: noem de twee delen apart. Extrema zijn er niet.

\end{alleenvoor}
\begin{oefeningen}
\needspace{6\baselineskip}
\oef{h2-homografisch}
\begin{alleenvoor}{2,3}

\needspace{6\baselineskip}
\oef{h2-k-homo}
\end{alleenvoor}
\end{oefeningen}

Omgekeerd kun je een functie ontwerpen met gegeven asymptoten: de vorm $c + \frac{a}{x - p}$ legt ze meteen vast, en een extra punt bepaalt $a$.

\voorbeeld{}{Geef een rationale functie met verticale asymptoot $x = -2$, horizontale asymptoot $y = 3$ en nulwaarde $1$.\label{vb:h2-ontwerp}}{%
\naast{%
\uitleg{Begin met een homografische functie. De asymptoten leggen $p$ en $c$ vast, de nulwaarde bepaalt $a$.}
\begin{stappenvrij}{5.4cm}{4.5cm}
f(x) = 3 + \frac{a}{x + 2} & $p = -2$, $c = 3$. \\[2ex]
f(1) = 3 + \frac{a}{3} = 0 \iff a = -9 & Nulwaarde $1$. \\[2ex]
f(x) = 3 - \frac{9}{x + 2} = \frac{3x - 3}{x + 2} & Op één noemer.
\end{stappenvrij}
Controle: de noemer is nul in $-2$, de teller niet; graden gelijk, $\frac31 = 3$; de teller is nul in $1$.}{%
\begin{tikzpicture}[x=.3cm,y=.3cm,font=\scriptsize]
  \rooster{-8}{2}{4}{-4}{2}{10}
  \draw[grijs,dashed] (-2,-4) -- (-2,10.5) (-8,3) -- (4.5,3);
  \kromme{-8}{-3.29}{3-9/(\x+2)} \kromme{-.71}{4}{3-9/(\x+2)}
  \vol{(1,0)}
\end{tikzpicture}}}

\begin{oefeningen}
\needspace{6\baselineskip}
\oef{h2-ontwerp}
\end{oefeningen}

\begin{alleenvoor}{2,3}
%========================================================================================
\section{Wortelfuncties}\label{sec:h2-wortel}

Een slinger met lengte $L$ meter zwaait in ongeveer $T(L) = 2\sqrt{L}$ seconden heen en terug. Een slinger van $1$ m doet er dus $2$ s over. Hoe lang moet hij zijn om er $4$ s over te doen? Niet $2$ m, want $T(2) = 2\sqrt2 \approx 2{,}83$ s. Pas bij $L = 4$ is $T = 4$: voor een slinger die twee keer zo traag zwaait, heb je een vier keer zo lange draad nodig. De functie stijgt, maar steeds trager. Een functie met de veranderlijke onder een wortel is geen rationale functie: ze heet \textit{irrationaal}. De eenvoudigste zijn de wortelfuncties.

\subsection{De grafiek van $\sqrt[n]{x}$}

Uit de cursus Exponentiële en logaritmische functies ken je de $n$-de machtswortel: $\sqrt[n]{x}$ is het getal $y$ met $y^n = x$, en bij een even $n$ neem je $y \geqslant 0$. Bij een even $n$ bestaat hij alleen voor $x \geqslant 0$, want een even macht is nooit negatief; bij een oneven $n$ bestaat hij voor elk reëel getal.

\begin{theorie}
De \textbf{\textcolor{accent}{$n$-de machtswortelfunctie}} is de functie \quad $\formule{f(x) = \sqrt[n]{x}}$ \quad met $n = 2, 3, 4, \ldots$
\end{theorie}

Omdat $y = \sqrt[n]{x}$ betekent dat $x = y^n$, is $\sqrt[n]{x}$ de inverse functie van $x^n$ (bij een even $n$ van $x^n$ met $x \geqslant 0$). De grafiek is dus het spiegelbeeld van die van $x^n$ om de rechte $y = x$. Zo ontstaat $\sqrt{x}$ uit de rechterhelft van de parabool $x^2$, en $\sqrt[3]{x}$ uit de hele grafiek van $x^3$.

\figuren{%
\begin{tikzpicture}[x=.55cm,y=.55cm,font=\scriptsize]
  \rooster{-1}{1}{4}{-1}{1}{4}
  \draw[grijs,densely dotted] (-1,-1) -- (4.2,4.2);
  \draw[grijs,thick,dashed] plot[domain=0:2,samples=40,smooth] (\x,{(\x)^2});
  \kromme{0}{4}{sqrt(\x)} \kromme[densely dashdotted]{0}{4}{(\x)^(1/4)}
  \node[accent,right] at (4,2) {$\sqrt{x}$};
  \node[accent,right] at (4,1.41) {$\sqrt[4]{x}$};
  \node[grijs,left] at (1.75,3.1) {$x^2$};
  \node[below] at (1.5,-1.3) {$n$ even};
\end{tikzpicture}\qquad\qquad
\begin{tikzpicture}[x=.55cm,y=.55cm,font=\scriptsize]
  \rooster{-3}{1}{3}{-2}{1}{2}
  \draw[grijs,densely dotted] (-2,-2) -- (2.4,2.4);
  \draw[grijs,thick,dashed] plot[domain=-1.26:1.26,samples=40,smooth] (\x,{(\x)^3});
  \kromme{0}{3}{(\x)^(1/3)} \kromme{-3}{0}{-((-\x)^(1/3))}
  \node[accent,right] at (3,1.44) {$\sqrt[3]{x}$};
  \node[grijs,left] at (1.2,1.9) {$x^3$};
  \node[below] at (0,-2.3) {$n$ oneven};
\end{tikzpicture}}

Op de figuren lees je de kenmerken af. Je ziet ook dat de grafiek steeds trager stijgt, en toch elke hoogte haalt.

\begin{theorie}
\begin{stelling}{Kenmerken van $\sqrt[n]{x}$}
Is $n$ \textbf{\textcolor{accent}{even}}, dan is \quad $\formule{\dom f = \ber f = \left[0, +\infty\right[}$

Is $n$ \textbf{\textcolor{accent}{oneven}}, dan is \quad $\formule{\dom f = \ber f = \R}$

In beide gevallen is $0$ de enige nulwaarde, is $f$ \textbf{\textcolor{accent}{stijgend}} op haar hele domein, en wordt $f(x)$ onbeperkt groot voor grote $x$.
\end{stelling}

\begin{bewijs}
\textit{Bereik.} Elk getal $y$ (bij een even $n$: elk $y \geqslant 0$) is een functiewaarde, want $\sqrt[n]{y^n} = y$. Bij een even $n$ is een wortel nooit negatief.

\textit{Stijgend.} Neem $x_1 < x_2$ in het domein, en $y_1 = \sqrt[n]{x_1}$, $y_2 = \sqrt[n]{x_2}$. Was $y_1 \geqslant y_2$, dan was ook $y_1^n \geqslant y_2^n$, want $x^n$ stijgt (bij een even $n$ voor $x \geqslant 0$). Dan stond er $x_1 \geqslant x_2$: dat kan niet. Dus is $y_1 < y_2$.

\textit{Onbeperkt groot.} Neem een hoogte $M > 0$. Voor $x > M^n$ is $\sqrt[n]{x} > \sqrt[n]{M^n} = M$, omdat $f$ stijgt.
\end{bewijs}
\end{theorie}

Een wortelfunctie heeft dus geen horizontale asymptoot. Ze stijgt wel steeds trager: $\sqrt{100} = 10$ en $\sqrt{10\,000} = 100$. Rechts van $0$ is ze afnemend stijgend. Ze heeft ook geen schuine asymptoot: voor grote $x$ blijft ze onder elke stijgende rechte, ook al is die nog zo vlak.

\subsection{Verschoven wortelfuncties}

Met de transformaties uit de cursus Algemene functies bouw je de grafiek van een verwante functie op, zoals bij de homografische functies.

\voorbeeld{}{Geef het domein, het bereik, de nulwaarde en het verloop van $f(x) = \sqrt{x + 4} - 2$, en schets de grafiek.\label{vb:h2-wortel}}{%
\naast{%
\uitleg{De grafiek is die van $\sqrt{x}$, $4$ naar links en $2$ naar onder geschoven. Onder een vierkantswortel mag niets negatiefs staan. Het beginpunt $(0, 0)$ van $\sqrt{x}$ schuift mee.}
\begin{stappenvrij}{4.4cm}{4.4cm}
x + 4 \geqslant 0 \iff x \geqslant -4 & $\dom f = \left[-4, +\infty\right[$. \\
(0, 0) \to (-4, -2) & Het beginpunt. \\
\ber f = \left[-2, +\infty\right[ & Want $\sqrt{x + 4} \geqslant 0$. \\
\sqrt{x + 4} = 2 \iff x + 4 = 4 & Kwadrateer: beide leden $\geqslant 0$. \\
\text{nulwaarde } 0 & Controle: $\sqrt4 - 2 = 0$. \\
\text{stijgend op } \left[-4, +\infty\right[ & Zoals $\sqrt{x}$.
\end{stappenvrij}}{%
\begin{tikzpicture}[x=.42cm,y=.42cm,font=\scriptsize]
  \rooster{-5}{1}{5}{-3}{1}{3}
  \draw[grijs,thick] plot[domain=0:5,samples=60,smooth] (\x,{sqrt(\x)});
  \kromme{-4}{5}{sqrt(\x+4)-2}
  \vol{(-4,-2)} \vol{(0,0)}
  \node[grijs,above] at (4,2.2) {$\sqrt{x}$};
\end{tikzpicture}}}

\begin{werkwijze}
Lees in $a\sqrt[n]{x - p} + c$ de transformaties af: verticaal uitrekken met $a$, $p$ naar rechts en $c$ naar boven. Bij een even $n$ volgt het domein uit $x - p \geqslant 0$, en is $(p, c)$ het beginpunt. Is $a > 0$, dan stijgt $f$; is $a < 0$, dan daalt ze. Voor de nulwaarde zet je de wortel alleen in één lid en verhef je tot de $n$-de macht. Controleer door in te vullen.
\end{werkwijze}

\begin{oefeningen}
\needspace{6\baselineskip}
\oef{h2-k-wortel}
\end{oefeningen}

%========================================================================================
\section{Grafisch oplossen met ICT}\label{sec:h2-ict}

Een drone brengt een bloedstaal naar een labo $4{,}8$ km verderop, en vliegt dan terug. Hij vliegt $15$ m/s ten opzichte van de lucht. Waait er een wind van $w$ m/s langs de route, dan vliegt hij heen aan $15 - w$ m/s en terug aan $15 + w$ m/s. Tijd is afstand gedeeld door snelheid, dus de vlucht duurt
\[t(w) = \frac{4800}{15 - w} + \frac{4800}{15 + w} \ \text{seconden} .\]
Bij windstilte is dat $640$ s. De piloot rekent met hoogstens $12$ minuten in de lucht. Tot welke windsnelheid lukt de vlucht? Je moet $t(w) \leqslant 720$ oplossen, met de onbekende in twee noemers. Daar heb je geen formule voor, maar wel een methode die je uit de cursus Algemene functies kent: teken beide leden als functie en zoek de snijpunten. Ze werkt voor elke vergelijking of ongelijkheid, ook met breuken en wortels.

\voorbeeld{}{Voor welke windsnelheden $w$, met $0 \leqslant w < 15$, duurt de vlucht van de drone hoogstens $720$ s? Los grafisch op met ICT.\label{vb:h2-ict}}{%
\naast{%
\uitleg{Teken het linkerlid en het rechterlid als functie. Kies een venster met $0 \leqslant w < 15$: bij een wind van $15$ m/s of meer komt de drone nooit aan. De ongelijkheid geldt waar de grafiek van $t$ onder de rechte ligt.}
\begin{stappenvrij}{4.2cm}{4cm}
y_1 = t(w), \quad y_2 = 720 & Beide leden als functie. \\
\text{snijpunt } (5, 720) & Op de TI-84: \texttt{intersect} in het menu CALC. \\
t(5) = 480 + 240 = 720 & Controle: hier exact. \\
0 \leqslant w \leqslant 5 & Daar ligt de grafiek op of onder de rechte.
\end{stappenvrij}}{%
\begin{tikzpicture}[x=.3cm,y=.2cm,font=\scriptsize]
  \draw[black!14,very thin,xstep=.3cm,ystep=.4cm] (0,0) grid (15,16);
  \draw[->,grijs] (0,0) -- (16.6,0) node[right,text=black] {$w$};
  \draw[->,grijs] (0,0) -- (0,17) node[above,text=black] {$t$};
  \foreach \t in {5,10,15}{\node[below,inner sep=1.5pt] at (\t,0) {$\t$};}
  \foreach \t in {4,8,12,16}{\node[left,inner sep=1.5pt] at (0,\t) {$\t00$};}
  \draw[grijs,dashed] (15,0) -- (15,16.5);
  \draw[accent,dashed] (0,7.2) -- (15,7.2);
  \kromme{0}{11.6}{1440/(225-(\x)^2)}
  \vol{(0,6.4)} \vol{(5,7.2)}
\end{tikzpicture}}
Bij een wind tot $5$ m/s ($18$ km/h) lukt de vlucht. Heen duurt ze dan $480$ s, terug $240$ s: de drone wint op de terugweg minder tijd dan hij heen verliest. Vergeet je het praktisch domein, dan zie je ook rechts van $15$ een stuk grafiek onder de rechte: daar is $t(w)$ negatief, en dat heeft geen betekenis.}

Let op bij een verticale asymptoot. Daar kan de ene grafiek van onder naar boven de andere springen zonder snijpunt, zoals bij $w = 15$. Sommige rekentoestellen verbinden de punten links en rechts van de asymptoot ook met een bijna verticale lijn; die hoort niet bij de grafiek.

\begin{werkwijze}
Teken beide leden als functie, in een venster dat de snijpunten en de verticale asymptoten toont. Zoek de snijpunten met ICT: hun $x$-coördinaten zijn de oplossingen van de vergelijking. Bij een ongelijkheid lees je af waar de ene grafiek boven of onder de andere ligt; let op de verticale asymptoten en op het domein. Controleer door in te vullen, en houd in een context alleen de oplossingen die betekenis hebben.
\end{werkwijze}

\begin{oefeningen}
\needspace{6\baselineskip}
\oef{h2-k-ict}
\end{oefeningen}

\end{alleenvoor}
%========================================================================================
\opdrachtenkop

\begin{oefeningen}
\needspace{14\baselineskip}
\oef{h2-o1}

\needspace{12\baselineskip}
\oef{h2-o10}

\needspace{16\baselineskip}
\oef{h2-o2}

\needspace{10\baselineskip}
\oef{h2-o12}

\needspace{6\baselineskip}
\oef{h2-o3}

\needspace{6\baselineskip}
\oef{h2-o11}

\begin{alleenvoor}{2,3}
\needspace{8\baselineskip}
\oef{h2-o32}

\end{alleenvoor}
\needspace{8\baselineskip}
\oef{h2-o9}

\needspace{12\baselineskip}
\oef{h2-o14}

\oef{h2-o4}

\begin{alleenvoor}{2,3}
\needspace{10\baselineskip}
\oef{h2-o28}

\end{alleenvoor}
\needspace{8\baselineskip}
\oef{h2-o19}

\begin{alleenvoor}{2,3}
\needspace{12\baselineskip}
\oef{h2-o33}

\end{alleenvoor}
\needspace{12\baselineskip}
\oef{h2-o17}

\oef{h2-o5}

\begin{alleenvoor}{2,3}
\needspace{8\baselineskip}
\oef{h2-o29}

\end{alleenvoor}
\needspace{6\baselineskip}
\oef{h2-o20}

\needspace{8\baselineskip}
\oef{h2-o21}

\needspace{6\baselineskip}
\oef{h2-o15}

\begin{alleenvoor}{2,3}
\needspace{10\baselineskip}
\oef{h2-o30}

\end{alleenvoor}
\needspace{8\baselineskip}
\oef{h2-o22}

\needspace{9\baselineskip}
\oef{h2-o18}

\oef{h2-o6}

\begin{alleenvoor}{3}
\needspace{10\baselineskip}
\oef{h2-o31}

\end{alleenvoor}
\needspace{12\baselineskip}
\oef{h2-o23}

\needspace{12\baselineskip}
\oef{h2-o16}

\oef{h2-o7}

\needspace{10\baselineskip}
\oef{h2-o25}

\needspace{10\baselineskip}
\oef{h2-o24}

\oef{h2-o27}

\needspace{8\baselineskip}
\oef{h2-o26}

\needspace{9\baselineskip}
\oef{h2-o8}
\end{oefeningen}

%----------------------------------------------------------------------------------------
% KOPPEL: verbanden met andere leerstof
%----------------------------------------------------------------------------------------
\needspace{20\baselineskip}
\opdrachtsectie{Koppel}

\lkr{Algemene functies, \alleenvoor{4}{Limieten en afgeleiden}\alleenvoor{2,3}{Afgeleiden} en Integralen: kijk na wat al gezien is. Het pakket legt de volgorde niet vast.}

\begin{oefeningen}
\oef{h2-koppel-1}

\oef{h2-koppel-2}

\oef{h2-koppel-3}
\end{oefeningen}

%----------------------------------------------------------------------------------------
\opdrachtsectie{Test jezelf}

\begin{oefeningen}
\oef{h2-t1}

\oef{h2-t2}

\oef{h2-t5}

\oef{h2-t3}

\oef{h2-t4}

\oef{h2-t6}
\begin{alleenvoor}{2,3}

\oef{h2-t7}

\oef{h2-t8}
\end{alleenvoor}
\end{oefeningen}

\herhaalschema{%
  \fv{oef:h2-t1}{vb:h2-perforatie}{oef:h2-domein}%
  \fv{oef:h2-t2}{vb:h2-ha}{oef:h2-ha}%
  \fv{oef:h2-t5}{vb:h2-homo}{oef:h2-homografisch}%
  \fv{oef:h2-t3}{vb:h2-sa}{oef:h2-sa}%
  \fv{oef:h2-t4}{vb:h2-schets}{oef:h2-schets}%
  \fv{oef:h2-t6}{vb:h2-ontwerp}{oef:h2-ontwerp}\alleenvoor{2,3}{\fv{oef:h2-t7}{vb:h2-wortel}{oef:h2-k-wortel}\fv{oef:h2-t8}{vb:h2-ict}{oef:h2-k-ict}}}

\ifzelfstudie
\section*{Maak je eigen samenvatting}
Sluit dit boek en schrijf in je schrift: wat een rationale functie is, hoe je domein, nulwaarden en perforaties bepaalt, de stellingen over verticale, horizontale en schuine asymptoten met hun bewijs, hoe je de kant van een verticale asymptoot vindt, waarom de deling de schuine asymptoot geeft, de zeven stappen om de grafiek te schetsen, en hoe je een homografische functie schrijft als $c + \frac{a}{x - p}$\alleenvoor{2,3}{ met haar bereik en waar ze stijgt of daalt, de kenmerken van $\sqrt[n]{x}$ met bewijs, hoe je een verschoven wortelfunctie schetst, en hoe je een vergelijking of ongelijkheid grafisch oplost met ICT}. Vergelijk daarna met de kaders.
\fi
